% Generated draft: requires mathematical and source review.
\documentclass[11pt]{article}
\usepackage[margin=25mm]{geometry}
\usepackage{fontspec}
\setmainfont{Latin Modern Roman}
\usepackage{amsmath,amssymb,mathtools,unicode-math}
\setmathfont{Latin Modern Math}
\usepackage{enumitem,needspace,xurl,hyperref}
\setlength{\parindent}{0pt}
\setlength{\parskip}{5pt}
\setlength{\emergencystretch}{3em}
\newcommand{\N}{\mathbb{N}}
\newcommand{\Z}{\mathbb{Z}}
\newcommand{\Q}{\mathbb{Q}}
\newcommand{\R}{\mathbb{R}}
\newcommand{\C}{\mathbb{C}}
\newcommand{\F}{\mathbb{F}}
\newcommand{\A}{\mathbb{A}}
\newcommand{\PP}{\mathbb{P}}
\newcommand{\E}{\mathbb{E}}
\newcommand{\eps}{\varepsilon}
\newcommand{\dd}{\,\mathrm{d}}
\begin{document}
\section*{UnsolvedMath ID: 800794}
\textbf{Source problem code:} catalogue entry
\subsection{Definitions and mathematical statement}
The supplied record defines simplicity explicitly but leaves the meanings of a Venn diagram and symmetry implicit. The following formulation uses a planar interpretation with finite curve intersections and rotational symmetry; these conventions require confirmation against the source.

Let \(n\geq 2\) be an integer. Consider distinct Jordan curves \(C_1,\ldots,C_n\subset\R^2\), with each pair meeting in finitely many points. Write \(D_i\) for the bounded component of \(\R^2\setminus C_i\). For each subset \(S\subseteq\{1,\ldots,n\}\), define
\[
U_S=\left\{x\in\R^2\setminus\bigcup_{i=1}^n C_i:\ x\in D_i\ \Longleftrightarrow\ i\in S\text{ for every }i\right\}.
\]
The family is a \emph{Venn diagram} if every \(U_S\) is nonempty and connected. It is \emph{simple} if
\[
\#\{i:x\in C_i\}\leq 2\qquad\text{for every }x\in\R^2.
\]
Here \emph{symmetric} means that, after relabeling, there is a point \(o\in\R^2\) such that the rotation \(\rho\) about \(o\) through angle \(2\pi/n\) satisfies
\[
\rho(C_i)=C_{i+1}\quad(1\leq i<n),\qquad \rho(C_n)=C_1.
\]
\textbf{Question (Kahle, 2019, Question 1).} For every prime integer \(n\), does there exist a simple, symmetric Venn diagram consisting of \(n\) curves, under these conventions?
\subsection{Short English statement}
Does every prime number \(n\) admit a planar Venn diagram with \(n\) curves, all \(2^n\) membership regions nonempty and connected, no point on more than two curves, and rotation through \(2\pi/n\) cyclically permuting the curves? The rotational interpretation and curve conventions require source confirmation.
\subsection{Sources}
\begin{itemize}
\item Local UnsolvedMath registry, supplied complete source record: permanent ID \(800794\), exact source problem code \texttt{catalogue entry}; upstream revision \texttt{d371ee917ac9647fbb06cd6fc529deec6f3dda82}. The registry attributes the question to Kahle (2019) and labels it open; this label is supplied metadata, not an independently verified assessment.
\item Kahle, \emph{Branko Gr\"unbaum in Many Dimensions: Open Problems} (2019), Question 1. Supplied source URL: \url{https://arxiv.org/abs/1901.08622}. The record reports access on 2026-10-16 and extraction by \texttt{arxiv\_latex}; these provenance claims are supplied and have not been independently verified.
\end{itemize}
\end{document}
